\ifdefined\gkver\else\def\gkver{student}\fi
\documentclass[\gkver]{gaokaozhenti}
\begin{document}

\chapter{2012年福建理综}
%% paperType: 理综物理部分

\begin{enumerate}
\item
%% number: 13
%% paperName: 2012年福建理综第 13 题 6 分
%% typeId: 1
%% type: 单选题
%% chapter: 机械波
%% point: 波的图象
%% method: 无
%% score: 6
%% degree: 650
%% duplicateId: 0
%% body:
一列简谐波沿$x$轴传播，$t=0$时刻的波形如图甲所示，此时质点P正沿$y$轴负方向运动，其振动图像如图乙所示，则该波的传播方向和波速分别是\xzanswer{A}
\onepicture{figs/13.png}
\fourchoices[answer=A]
{沿$x$轴负方向，$60\Ums$}
{沿$x$轴正方向，$60\Ums$}
{沿$x$轴负方向，$30\Ums$}
{沿$x$轴正方向，$30\Ums$}
%% answer: A
%% memo:
\memoanswer{
根据波的形成和传播规律可知，波沿$x$轴负向传播，B、D错误；由图甲可知波长$\lambda=24\Um$，由图乙可知周期$T=(0.55-0.15)\Us=0.40\Us$，则波速$v=\frac{\lambda}{T}=\frac{24}{0.40}\Ums=60\Ums$，选项A正确，C错。
}

\item
%% number: 14
%% paperName: 2012年福建理综第 14 题 6 分
%% typeId: 1
%% type: 单选题
%% chapter: 交流电
%% point: 变压器
%% method: 无
%% score: 6
%% degree: 650
%% duplicateId: 0
%% body:
如图，理想变压器原线圈输入电压$u=U_{m}\sin\omega t$，副线圈电路中$R_{0}$为定值电阻，$R$是滑动变阻器，$V_{1}$和$V_{2}$是理想交流电压表，示数分别用$U_{1}$和$U_{2}$表示；$A_{1}$和$A_{2}$是理想交流电流表，示数分别用$I_{1}$和$I_{2}$表示。下列说法正确的是\xzanswer{C}
\onepicture{figs/14.png}
\fourchoices[answer=C]
{$I_{1}$和$I_{2}$表示电流的瞬时值}
{$U_{1}$和$U_{2}$表示电压的最大值}
{滑片$P$向下滑动过程中，$U_{2}$不变、$I_{1}$变大}
{滑片$P$向下滑动过程中，$U_{2}$不变、$I_{1}$变小}
%% answer: C
%% memo:
\memoanswer{
交流电表的示数表示交流电的有效值，选项A、B错误；滑片P向下滑动过程中，原、副线圈的匝数$n_{1}$、$n_{2}$不变，$U_{1}$不变，则$U_{2}=\frac{n_{2}}{n_{1}}U_{1}$不变；R连入电路的电阻减小，则流过副线圈的电流$I_{2}=\frac{U_{2}}{R+R_{0}}$变大，根据$U_{1}I_{1}=U_{2}I_{2}$，可知$I_{1}$变大，选项C正确，D错。
}

\item
%% number: 15
%% paperName: 2012年福建理综第 15 题 6 分
%% typeId: 1
%% type: 单选题
%% chapter: 电场
%% point: 等势面
%% method: 无
%% score: 6
%% degree: 650
%% duplicateId: 0
%% body:
如图，在点电荷$Q$产生的电场中，将两个带正电的试探电荷$q_{1}$、$q_{2}$分别置于A、B两点，虚线为等势线。取无穷远处为零电势点，若将$q_{1}$、$q_{2}$移动到无穷远的过程中外力克服电场力做的功相等，则下列说法正确的是\xzanswer{C}
\onepicture{\includesvg{figs/15}}
\fourchoices[answer=C]
{A点电势大于B点电势}
{A、B两点的电场强度相等}
{$q_{1}$的电荷量小于$q_{2}$的电荷量}
{$q_{1}$在A点的电势能小于$q_{2}$在B点的电势能}
%% answer: C
%% memo:
\memoanswer{
由外力克服电场力做功可知电场力做负功，点电荷Q带负电，电场线指向Q，根据“沿着电场线电势降低”可知B点的电势高于A点的电势，选项A错误；根据场强$E=\frac{kQ}{r^2}$可知，距离Q较近的A点的场强较大，选项B错误；电荷在某点的电势能等于把电荷从该点移到零电势能点的过程中电场力做的功，可知两个电荷在两点的电势能相等，选项D错误。根据电势$\varphi=\frac{E_p}{q}$，$\varphi_B>\varphi_A$可知，$q_1$的电荷量小于$q_2$的电荷量，选项C正确。
}

\item
%% number: 16
%% paperName: 2012年福建理综第 16 题 6 分
%% typeId: 1
%% type: 单选题
%% chapter: 曲线运动
%% point: 人造地球卫星
%% method: 无
%% score: 6
%% degree: 450
%% duplicateId: 0
%% body:
一卫星绕某一行星表面附近做匀速圆周运动，其线速度大小为$v$。假设宇航员在该行星表面上用弹簧测力计测量一质量为$m$的物体重力，物体静止时，弹簧测力计的示数为$N$。已知引力常量为$G$，则这颗行星的质量为\xzanswer{B}
\fourchoices[answer=B]
{$\frac{mv^2}{GN}$}
{$\frac{mv^4}{GN}$}
{$\frac{Nv^2}{Gm}$}
{$\frac{Nv^4}{Gm}$}
%% answer: B
%% memo:
\memoanswer{
行星对卫星的万有引力提供其做匀速圆周运动的向心力，有$G\frac{Mm'}{R^2}=m'\frac{v^2}{R}$；行星对处于其表面物体的万有引力等于物体重力，有$G\frac{Mm}{R^2}=mg$；而$N=mg$，可得$M=\frac{mv^4}{GN}$，选项B正确。
}

\item
%% number: 17
%% paperName: 2012年福建理综第 17 题 6 分
%% typeId: 1
%% type: 单选题
%% chapter: 功和能
%% point: 功率
%% method: 无
%% score: 6
%% degree: 450
%% duplicateId: 0
%% body:
如图，表面光滑的固定斜面顶端安装一定滑轮，小物块$A$、$B$用轻绳连接并跨过滑轮（不计滑轮的质量和摩擦）。初始时刻，$A$、$B$处于同一高度并恰好处于静止状态。剪断轻绳后$A$下落、$B$沿斜面下滑，则从剪断轻绳到物块着地，两物块\xzanswer{D}
\onepicture{figs/17.png}
\fourchoices[answer=D]
{速率的变化量不同}
{机械能的变化量不同}
{重力势能的变化量相同}
{重力做功的平均功率相同}
%% answer: D
%% memo:
\memoanswer{
A、B静止有$m_Ag=m_Bg\sin\theta$，二者速率的变化量（加速度大小）分别为$g$、$g\sin\theta$，它们相等，选项A错误；两个物体下滑过程中机械能分别守恒，因而机械能的变化量都为零，选项B错误；重力势能的变化量等于重力做的功，分别为$W_A=m_Agh$、$W_B=m_Bgh$，又$m_A<m_B$，后者较大，选项C错误；由$h=\frac{1}{2}gt_A^2$、$\frac{h}{\sin\theta}=\frac{1}{2}g\sin\theta t_B^2$可得A、B下滑时间；根据平均功率$P=W/t$可得$P_A=P_B$，选项D正确。
}

\item
%% number: 18
%% paperName: 2012年福建理综第 18 题 6 分
%% typeId: 1
%% type: 单选题
%% chapter: 电磁感应
%% point: 电磁感应中图象
%% method: 无
%% score: 6
%% degree: 450
%% duplicateId: 0
%% body:
如图甲，一圆形闭合铜环由高处从静止开始下落，穿过一根竖直悬挂的条形磁铁，铜环的中心轴线与条形磁铁的中轴线始终保持重合。若取磁铁中心O为坐标原点，建立竖直向下为正方向的$x$轴，则图乙中最能正确反映环中感应电流$i$随环心位置坐标$x$变化的关系图象是\xzanswer{B}
\onepicture{figs/18.png}
\fourchoices[answer=B, ispicture=true, h=2.0cm]
{figs/18a.png}
{figs/18b.png}
{figs/18c.png}
{figs/18d.png}
%% answer: B
%% memo:
\memoanswer{
圆环下落过程中，穿过的磁通量先增大后减小，电流方向为先顺时针后逆时针（从上往下看），选项D错误。圆环通过O位置时，不切割磁感线，没有感应电流，只受重力作用，而在关于O对称的位置上穿过圆环的磁通量相等，磁通量的变化率随圆环速度的不同而不同。在正$x$轴上的相应位置上，圆环的速度较大，其磁通量的变化率较大，因而感应电流的最大值大于圆环在负$x$轴上的感应电流最大值，选项B正确，A、C错误。
}

\item
%% number: 19
%% paperName: 2012年福建理综第 19 题 6 分
%% typeId: 4
%% type: 实验题
%% chapter: 光学
%% point: 双缝干涉测量波长
%% method: 无
%% score: 6
%% degree: 650
%% duplicateId: 0
%% body:
在“用双缝干涉测光的波长”实验中：（实验装置如下图）
\onepicture[h=2.2cm]{figs/19a.png}
\begin{enumerate}
	\item
	下列说法哪一个是错误的\underline{\hspace{1.6cm}}。（填选项前的字母）

	（A）调节光源高度使光束沿遮光筒轴线照在屏中心时，应放上单缝和双缝

	（B）测量某条干涉亮纹位置时，应使测微目镜分划板中心刻线与该亮纹的中心对齐

	（C）为了减小测量误差，可用测微目镜测出$n$条亮纹间的距离$a$，求出相邻两条亮纹间距$\Delta x=\frac{a}{n-1}$
	\item
	测量某亮纹位置时，手轮上的示数如图，其示数为\underline{\hspace{1.6cm}}$\Umm$。
\end{enumerate}
\onepicture[h=2.8cm, align=right]{figs/19b.png}
%% answer: 
\jdanswer{
\begin{enumerate}
	\item
	A
	\item
	$1.970\Umm$
\end{enumerate}
}
%% memo:
\memoanswer{
①调节光源高度使光束沿遮光筒轴线照在屏中心时，无须放上单缝和双缝。②主尺的示数为$1.5\Umm$（半毫米刻度线已经露出），可动尺的示数为$47.0\times0.01\Umm=0.470\Umm$，总的示数为$(1.5+0.470)\Umm=1.970\Umm$。
}

\item
%% number: 19
%% paperName: 2012年福建理综第 19 题 12 分
%% typeId: 4
%% type: 实验题
%% chapter: 电路
%% point: 测电动势与内阻
%% method: 无
%% score: 12
%% degree: 450
%% duplicateId: 0
%% body:
某研究性学习小组欲测定一块电池的电动势$E$。

（1）先直接用多用电表测定该电池电动势。在操作无误的情况下，多用电表表盘示数如图，其示数为\underline{\hspace{1.6cm}}V。
\onepicture[h=4.0cm, align=right]{figs/19d.png}
（2）然后，用电压表V、电阻箱$R$、定值电阻$R_{0}$、开关S、若干导线和该电池组成电路，测定该电池电动势。

（i）根据电路图，用笔画线代替导线，将实物图连接成完整电路。
\onepicture[h=2.4cm]{figs/19e.png}
（ii）闭合开关S，调整电阻箱阻值$R$，读出电压表V相应示数$U$。该学习小组测出大量数据，分析筛选出下表所示的$R$、$U$数据，并计算出相应的$\frac{1}{R}$与$\frac{1}{U}$的值。请用表中数据在坐标纸上描点，并作出$\frac{1}{U}-\frac{1}{R}$图线。

\begin{center}
\begin{tabular}{|c|c|c|c|c|c|c|}
\hline
$R$（$\Omega$） & 166.7 & 71.4 & 50.0 & 33.3 & 25.0 & 20.0 \\
\hline
$U$（V） & 8.3 & 5.9 & 4.8 & 4.2 & 3.2 & 2.9 \\
\hline
$\frac{1}{R}$（$\times10^{-2}\Omega^{-1}$） & 0.60 & 1.40 & 2.00 & 3.00 & 4.00 & 5.00 \\
\hline
$\frac{1}{U}$（V$^{-1}$） & 0.12 & 0.17 & 0.21 & 0.24 & 0.31 & 0.35 \\
\hline
\end{tabular}
\end{center}
\onepicture[h=3.4cm]{figs/19c.png}
（iii）从图线中可求得$E=$\underline{\hspace{1.6cm}}V。
%% answer: 
\jdanswer{
①$9.4\UV$；②（i）见答图1；（ii）见答图2；（iii）$9.5\sim11.1\UV$。
}
%% memo:
\memoanswer{
①选择开关位于直流电压$10\UV$档，按电表中央刻度$0\sim10\UV$刻度读数，最小刻度为$0.2\UV$，电池的电动势为$9.4\UV$。

②（i）连接电路如图所示。
\onepicture[h=2.4cm, align=right]{figs/19_an1.png}

（ii）图象如图所示。
\onepicture[h=3.2cm, align=right]{figs/19_an2.png}

（iii）根据闭合电路欧姆定律有$E=U+\frac{U}{R}R_0$，化简得$\frac{1}{U}=\frac{1}{E}+\frac{R_0}{E}\cdot\frac{1}{R}$，可知$\frac{1}{U}-\frac{1}{R}$图线在纵轴的截距$b=1/E$。由图线可知$b=0.10$，则$E=10\UV$。
}

\item
%% number: 20
%% paperName: 2012年福建理综第 20 题 15 分
%% typeId: 6
%% type: 计算题
%% chapter: 曲线运动
%% point: 平抛运动
%% method: 无
%% score: 15
%% degree: 650
%% duplicateId: 0
%% body:
如图，置于圆形水平转台边缘的小物块随转台加速转动，当转速达到某一数值时，物块恰好滑离转台开始做平抛运动。现测得转台半径$R=0.5\Um$，离水平地面的高度$H=0.8\Um$，物块平抛落地过程水平位移的大小$s=0.4\Um$。设物体所受的最大静摩擦力等于滑动摩擦力，取重力加速度$g=10\Umsq$。求：
\begin{enumerate}
	\item
	物块做平抛运动的初速度大小$v_{0}$；
	\item
	物块与转台间的动摩擦因数$\mu$。
\end{enumerate}
\onepicture[align=right]{figs/20.png}
%% answer: 
\jdanswer{
\begin{enumerate}
	\item
	$1\Ums$
	\item
	$0.2$
\end{enumerate}
}
%% memo:
\memoanswer{
（1）物块做平抛运动，在竖直方向上有$H=\frac{1}{2}gt^{2}$，在水平方向上有$s=v_{0}t$，解得$v_{0}=s\sqrt{\frac{g}{2H}}=1\Ums$。

（2）物块离开转台时，最大静摩擦力提供向心力，有$F_{m}=m\frac{v_{0}^{2}}{R}$，又$F_{m}=F_{f}=\mu F_{N}=\mu mg$，解得$\mu=\frac{v_{0}^{2}}{gR}=0.2$。
}

\item
%% number: 21
%% paperName: 2012年福建理综第 21 题 19 分
%% typeId: 6
%% type: 计算题
%% chapter: 功和能
%% point: 动能定理
%% method: 无
%% score: 19
%% degree: 450
%% duplicateId: 0
%% body:
如图，用跨过光滑定滑轮的缆绳将海面上一艘失去动力的小船沿直线拖向岸边。已知拖动缆绳的电动机功率恒为$P$，小船的质量为$m$，小船受到的阻力大小恒为$f$，经过A点时的速度大小为$v_{0}$，小船从A点沿直线加速运动到B点经历时间为$t_{1}$，A、B两点间距离为$d$，缆绳质量忽略不计。求：
\begin{enumerate}
	\item
	小船从A点运动到B点的全过程克服阻力做的功$W_{f}$；
	\item
	小船经过B点时的速度大小$v_{1}$；
	\item
	小船经过B点时的加速度大小$a$。
\end{enumerate}
\onepicture[align=right]{figs/21.png}
%% answer: 
\jdanswer{
\begin{enumerate}
	\item
	$fd$
	\item
	$\sqrt{v_0^2+\frac{2}{m}(Pt_1-fd)}$
	\item
	$\frac{P}{\sqrt{m^2v_0^2+2m(Pt_1-fd)}}-\frac{f}{m}$
\end{enumerate}
}
%% memo:
\memoanswer{
（1）小船从A点运动到B点克服阻力做功$W_{f}=fd$。

（2）小船从A点运动到B点，电动机牵引绳对小船做功$W=Pt_{1}$，由动能定理有$W-W_{f}=\frac{1}{2}mv_{1}^{2}-\frac{1}{2}mv_{0}^{2}$，解得$v_{1}=\sqrt{v_{0}^{2}+\frac{2}{m}(Pt_{1}-fd)}$。

（3）设小船经过B点时绳的拉力大小为$F$，绳与水平方向夹角为$\theta$，电动机牵引绳的速度大小为$u$，则$P=Fu$，$u=v_{1}\cos\theta$；由牛顿第二定律有$F\cos\theta-f=ma$，由以上各式得$a=\frac{P}{\sqrt{m^{2}v_{0}^{2}+2m(Pt_{1}-fd)}}-\frac{f}{m}$。
}

\item
%% number: 22
%% paperName: 2012年福建理综第 22 题 20 分
%% typeId: 6
%% type: 计算题
%% chapter: 磁场
%% point: 带电粒子在磁场中的运动
%% method: 无
%% score: 20
%% degree: 450
%% duplicateId: 0
%% body:
如图甲，在圆柱形区域内存在一方向竖直向下、磁感应强度大小为$B$的匀强磁场，在此区域内，沿水平面固定一半径为$r$的圆环形光滑细玻璃管，环心O在区域中心。一质量为$m$、带电荷量为$q$（$q>0$）的小球，在管内沿逆时针方向（从上向下看）做圆周运动。已知磁感应强度大小$B$随时间$t$的变化关系如图乙所示，其中$T_{0}=\frac{2\pi m}{qB_{0}}$。设小球在运动过程中电荷量保持不变，对原磁场的影响可忽略。
\begin{enumerate}
	\item
	在$t=0$到$t=T_{0}$这段时间内，小球不受细管侧壁的作用力，求小球的速度大小$v_{0}$；
	\item
	在竖直向下的磁感应强度增大过程中，将产生涡旋电场，其电场线是在水平面内一系列沿逆时针方向的同心圆，同一条电场线上各点的场强大小相等。试求$t=T_{0}$到$t=1.5T_{0}$这段时间内：
	\begin{enumerate}
		\item
		细管内涡旋电场的场强大小$E$；
		\item
		电场力对小球做的功$W$。
	\end{enumerate}
\end{enumerate}
\onepicture[align=right]{figs/22.png}
%% answer: 
\jdanswer{
\begin{enumerate}
	\item
	$\frac{qB_0r}{m}$
	\item
	①$\frac{qB_0^2r}{2\pi m}$；②$\frac{5q^2B_0^2r^2}{8m}$
\end{enumerate}
}
%% memo:
\memoanswer{
（1）小球运动时不受细管侧壁的作用力，因而小球所受洛伦兹力提供向心力$qv_{0}B_{0}=\frac{mv_{0}^{2}}{r}$，由上式得$v_{0}=\frac{qB_{0}r}{m}$。

（2）①在$T_{0}$到$1.5T_{0}$这段时间内，细管内一周的感应电动势$E_{\text{感}}=\pi r^{2}\frac{\Delta B}{\Delta t}$，由图乙可知$\frac{\Delta B}{\Delta t}=\frac{2B_{0}}{T_{0}}$；由于同一电场线上各点的场强大小相等，所以$E=\frac{E_{\text{感}}}{2\pi r}$，由以上各式及$T_{0}=\frac{2\pi m}{B_{0}q}$得$E=\frac{qB_{0}^{2}r}{2\pi m}$。

②在$T_{0}$到$1.5T_{0}$这段时间内，小球沿切线方向的加速度大小恒为$a=\frac{Eq}{m}$，小球运动的末速度大小$v=v_{0}+a\Delta t$，由图乙$\Delta t=0.5T_{0}$，得$v=\frac{3}{2}v_{0}=\frac{3qB_{0}r}{2m}$；由动能定理，电场力做功为$W=\frac{1}{2}mv^{2}-\frac{1}{2}mv_{0}^{2}$，由以上各式解得$W=\frac{5}{8}mv_{0}^{2}=\frac{5q^{2}B_{0}^{2}r^{2}}{8m}$。
}

\item
%% number: 28
%% paperName: 2012年福建理综第 28 题 5 分
%% typeId: 1
%% type: 单选题
%% chapter: 气体性质
%% point: 分子力
%% method: 无
%% score: 5
%% degree: 650
%% duplicateId: 0
%% body:
关于热力学定律和分子动理论，下列说法正确的是\tkanswer{D}。（填选项前的字母）
\fourchoices
{一定量气体吸收热量，其内能一定增大}
{不可能使热量由低温物体传递到高温物体}
{若两分子间距离增大，分子势能一定增大}
{若两分子间距离减小，分子间引力和斥力都增大}
%% answer : D
%% memo:
\memoanswer{
一定量气体吸收热量的同时对外做功，则其内能可能不变或减小，选项A错误；在外界做功的前提下，热量可以由低温物体传递到高温物体，选项B错误；若两分子间距增大，分子力可能做正功或负功，其势能可能减小或增大，选项C错误。根据分子动理论可知，若分子之间的距离减小，分子间的引力和斥力都增大，选项D正确。
}

\item
%% number: 28
%% paperName: 2012年福建理综第 28 题 5 分
%% typeId: 1
%% type: 单选题
%% chapter: 气体性质
%% point: 玻意耳定律
%% method: 无
%% score: 5
%% degree: 450
%% duplicateId: 0
%% body:
空气压缩机的储气罐中储有$1.0\Uatm$的空气$6.0\UL$，现再充入$1.0\Uatm$的空气$9.0\UL$。设充气过程为等温过程，空气可看作理想气体，则充气后储气罐中气体压强为\tkanswer{A}。（填选项前的字母）
\fourchoices
{$2.5\Uatm$}
{$2.0\Uatm$}
{$1.5\Uatm$}
{$1.0\Uatm$}
%% answer : A
%% memo:
\memoanswer{
根据玻意耳定律有$P_{1}V_{1}+P_{2}V_{2}=P_{3}V_{3}$，即$1\times6+1\times9=6P_{3}$，解得$P_{3}=2.5\Uatm$。
}

\item
%% number: 29
%% paperName: 2012年福建理综第 29 题 5 分
%% typeId: 1
%% type: 单选题
%% chapter: 原子和原子核
%% point: 核裂变
%% method: 无
%% score: 5
%% degree: 650
%% duplicateId: 0
%% body:
关于近代物理，下列说法正确的是\tkanswer{D}。（填选项前的字母）
\fourchoices
{$\alpha$射线是高速运动的氦原子}
{核聚变反应方程$\ce{^{2}_{1}H + ^{3}_{1}H -> ^{4}_{2}He + ^{1}_{0}n}$中，$\ce{^{1}_{0}n}$表示质子}
{从金属表面逸出的光电子的最大初动能与照射光的频率成正比}
{玻尔将量子观念引入原子领域，其理论能够解释氢原子光谱的特征}
%% answer : D
%% memo:
\memoanswer{
$\alpha$射线是高速运动的氦原子核，选项A错误；选项B中$\ce{^{1}_{0}n}$表示中子；由光电方程$\frac{1}{2}mv_{m}^{2}=h\nu-W$可知最大初动能与照射光的频率成线性关系而非正比，选项C错误；根据玻尔原子理论可知，选项D正确。
}

\item
%% number: 29
%% paperName: 2012年福建理综第 29 题 5 分
%% typeId: 1
%% type: 单选题
%% chapter: 运动和力
%% point: 动量守恒定律
%% method: 无
%% score: 5
%% degree: 450
%% duplicateId: 0
%% body:
如图，质量为$M$的小船在静止水面上以速率$v_{0}$向右匀速行驶，一质量为$m$的救生员站在船尾，相对小船静止。若救生员以相对水面速率$v$水平向左跃入水中，则救生员跃出后小船的速率为\tkanswer{C}。（填选项前的字母）
\onepicture{figs/29.png}
\fourchoices
{$v_{0}+\frac{m}{M}v$}
{$v_{0}-\frac{m}{M}v$}
{$v_{0}+\frac{m}{M}(v_{0}+v)$}
{$v_{0}+\frac{m}{M}(v_{0}-v)$}
%% answer : C
%% memo:
\memoanswer{
设水平向右为正方向，根据动量守恒定律，对救生员和船有$(M+m)v_{0}=-mv+Mv_{x}$，解得$v_{x}=v_{0}+\frac{m}{M}(v_{0}+v)$。选项C正确。
}

\end{enumerate}

\end{document}
